\documentclass[12pt]{article}
\usepackage[margin=1in]{geometry}
\usepackage{booktabs,graphicx,threeparttable}

\begin{document}

\providecommand{\StructuralRobustnessPath}{appendix/structural-robustness}

\subsection{What Do the Signals Mean? Coding the Free-Text Meaning Responses}
\label{app:signal-meaning-coding}

At endline, every respondent who reported having seen a bracelet, ink mark, or
calendar was asked what the item means. This yields 2{,}239 free-text answers
(950 bracelet, 620 ink, 669 calendar), each attributable to a community whose
distance condition was randomized. We use these responses to ask two questions
that bear directly on the interpretation of the structural decomposition.
First, does the meaning respondents attach to a public signal include
attributions about the participant's motivation or character, and does that
content strengthen with randomized distance? Second, is the basic denotation
of the signals stable across distance conditions?

The coding is designed to be auditable. All 831 distinct response strings were
extracted into a coding sheet containing only an identifier, the item, and the
verbatim text---no treatment, distance, or community fields---so coders could
not condition on the experimental assignment. A codebook fixed before any
coding defines five non-exclusive categories: motivation or character
attribution (M), participation marker (P), campaign or verification function
(C), private or decorative item value (V), and other or uncodable (O), with a
prespecified conservative rule that withholds M in borderline cases. Two
independent language-model coders labeled every string from the codebook
alone; they disagreed on 5 of 831 strings (exact five-label agreement 99.4
percent; per-label agreement at least 99.6 percent). A third blind pass
adjudicated the disagreements by majority vote. The coded artifacts are
committed with the replication materials, so the analysis reproduces without
re-running any coder.

Table~\ref{tab:signal-meaning-shares} reports the category shares by item and
original distance assignment; Table~\ref{tab:signal-meaning-quotes} gives the
most common examples. The result is unambiguous: the meaning of every item is
overwhelmingly a bare participation marker. Between 89 and 98 percent of
responses in every item-by-distance cell state that the holder was dewormed,
with small verification (C) and gift or reminder (V) components---the latter
concentrated in the calendar arm, consistent with its private-value
interpretation. Spontaneous motivation or character attributions are nearly
absent: 3 of 831 distinct strings, or well under one percent of respondents in
every cell.

\begin{table}[htbp]
  \centering
  \small
  \caption{Coded meaning content by item and original distance assignment}
  \label{tab:signal-meaning-shares}
  \begin{threeparttable}
    \begin{tabular}{llcccccc}
\toprule
Item & Assignment & $N$ & M & P & C & V & O \\
\midrule
Bracelet & Close & 470 & 0.000 & 0.974 & 0.011 & 0.013 & 0.019 \\
 & Far & 480 & 0.004 & 0.952 & 0.010 & 0.013 & 0.027 \\
\addlinespace
Ink & Close & 411 & 0.000 & 0.956 & 0.027 & 0.000 & 0.032 \\
 & Far & 209 & 0.000 & 0.981 & 0.010 & 0.000 & 0.014 \\
\addlinespace
Calendar & Close & 382 & 0.000 & 0.887 & 0.005 & 0.123 & 0.029 \\
 & Far & 287 & 0.003 & 0.909 & 0.000 & 0.073 & 0.028 \\
\bottomrule
\end{tabular}

    \begin{tablenotes}[flushleft]\footnotesize
      \item \textit{Notes:} Entries are respondent-level shares of each coded
      category among endline respondents who saw the item, by original
      randomized Close/Far assignment. Categories are non-exclusive except O.
      M: motivation or character attribution; P: participation marker; C:
      campaign or verification function; V: private, decorative, or reminder
      value; O: other or uncodable.
    \end{tablenotes}
  \end{threeparttable}
\end{table}

\begin{table}[htbp]
  \centering
  \small
  \caption{Distance contrasts in coded meaning content and item exposure}
  \label{tab:signal-meaning-effects}
  \resizebox{\linewidth}{!}{%
    \begin{tabular}{lccccccc}
\toprule
Sample & Outcome & Close mean & Far $-$ Close & 95\% interval & Fisher $p$ & Stud.\ RI $p$ & $N$ \\
\midrule
Bracelet (primary) & Motivation (M) & 0.000 & 0.003 & [0.000, 0.009] & 0.368 & 0.265 & 727 \\
Bracelet, incl. SMS & Motivation (M) & 0.000 & 0.005 & [0.000, 0.015] & 0.297 & 0.219 & 950 \\
Bracelet, realized Far & Motivation (M) & 0.000 & 0.003 & [0.000, 0.011] & -- & -- & 727 \\
Ink & Motivation (M) & 0.000 & 0.000 & [0.000, 0.000] & 1.000 & -- & 445 \\
Pooled signals & Motivation (M) & 0.000 & 0.002 & [0.000, 0.006] & 0.547 & 0.414 & 1,172 \\
Calendar (placebo) & Motivation (M) & 0.000 & 0.005 & [0.000, 0.016] & 0.059 & 0.073 & 497 \\
Calendar, private value & Private value (V) & 0.111 & -0.035 & [-0.103, 0.025] & 0.396 & 0.352 & 497 \\
Bracelet, marker only & Marker only & 0.961 & -0.022 & [-0.051, 0.006] & 0.209 & 0.229 & 727 \\
Bracelet, unconditional & Motivation (M) & 0.000 & 0.003 & [0.000, 0.009] & 0.409 & 0.386 & 769 \\
\addlinespace
Seen bracelet & Seen item & 0.960 & -0.037 & [-0.068, -0.003] & -- & -- & 769 \\
Seen ink & Seen item & 0.696 & -0.123 & [-0.216, -0.029] & -- & -- & 677 \\
Seen calendar & Seen item & 0.847 & -0.229 & [-0.313, -0.135] & -- & -- & 684 \\
\bottomrule
\end{tabular}
}
  \begin{minipage}{0.96\linewidth}
    \vspace{0.35em}\footnotesize
    \textit{Notes:} Far $-$ Close contrasts under the original randomized
    assignment unless labeled otherwise, adjusting for sex, age, expected
    distance, and county fixed effects. Intervals are 95 percent percentile
    intervals across 1{,}000 county-stratified exponential community-weight
    draws. Randomization-inference $p$-values permute the Close/Far label
    across the sample's communities within county; at the near-zero base rate
    of M they are driven by one or two responses and are reported only for
    completeness. The primary sample excludes SMS-treated respondents. ``Seen
    item'' rows use all arm respondents and report whether the respondent
    reported having seen the item.
  \end{minipage}
\end{table}

\begin{table}[htbp]
  \centering
  \small
  \caption{Most common coded examples (bracelet arm)}
  \label{tab:signal-meaning-quotes}
  \begin{tabular}{lp{0.62\linewidth}r}
\toprule
Category & Example responses (verbatim) & $N$ \\
\midrule
M & ``He's taken care of his Health'' & 1 \\
 & ``One is taking care of his Health'' & 1 \\
\addlinespace
P & ``Dewormed'' & 53 \\
 & ``They have been dewormed'' & 49 \\
 & ``They were dewormed'' & 49 \\
\addlinespace
C & ``Confirmation that someone has gone to take deworming tablet'' & 1 \\
 & ``Has dewormed. Emphasizes deworming'' & 1 \\
 & ``Has received dewormimg Advertising there is deworming treatment'' & 1 \\
\addlinespace
V & ``A preent;\& a sign that dewormed'' & 1 \\
 & ``Aesthetic purposes'' & 1 \\
 & ``Appreciation as after taking drug'' & 1 \\
\addlinespace
O & ``Don't know'' & 8 \\
 & ``I don't know'' & 2 \\
 & ``1'' & 1 \\
\bottomrule
\end{tabular}

\end{table}

Table~\ref{tab:signal-meaning-effects} reports the experimental contrasts.
Three findings follow. First, the near-zero motivation share does not differ
by randomized distance for the bracelet (0.3 percentage points, interval
$[0.0, 0.9]$), for ink, for the pooled signals, or for the calendar placebo;
at this base rate the exercise cannot detect gradients, and none is visible.
Second, the denotation of the signals is stable: the participation-marker
share is 95 to 98 percent in every signal cell, Close and Far alike. The
bracelet does not become a virtue badge at distance---or anywhere. This
stability supports the structural model's maintained assumption that the
public signals change the visibility of participation rather than the meaning
of participating, and complements the signal-specific image-weight analysis in
Section~\ref{app:signal-specific-lambda}. Third, reported exposure to the
items declines with randomized distance (by 3.7 percentage points for the
bracelet, 12.3 for ink, and 22.9 for the calendar), an independent,
non-model-based corroboration that distance reduces the circulation of
participation information.

Two limitations define what this exercise can and cannot establish. The
question elicits the signal's denotation, not an evaluation of its wearer;
respondents who would judge a Far participant more favorably need not
volunteer that judgment when asked what a bracelet means. And the contrasts
condition on having seen the item, which itself responds to distance; the
unconditional bracelet estimate, which codes non-seers as zero, is reported in
Table~\ref{tab:signal-meaning-effects} and is unchanged. We therefore draw a
deliberately limited conclusion: the free-text evidence shows no direct trace
of distance-dependent motivation attribution, and the type-inference component
of the multiplier should be read as a model-implied decomposition---while the
meaning stability and exposure gradients independently support the
observability channel on which that decomposition rests.

\end{document}
